% id: 91-01
% book: 베이직쎈 미적분1 · 05 도함수의 활용 (2)
% section: 학교 시험 기출로 실전 감각 UP!
% figure: tikz:fig-91-01
% answer: 
% answer_source: 
% variant_level: 0 (원본 전사)
% 이 파일은 build.mjs 가 items.json 에서 만든다. 고칠 때는 items.json 을 고친다.
\smash{\raisebox{1.5mm}{\dmkichul{베이직쎈 미적분1 91쪽 01번}}}
\noindent\begin{minipage}[t]{0.58\linewidth}\vspace{0pt}\raggedright 다항함수 $f(x)$의 도함수 $y=f'(x)$의 그래프가 오른쪽 그림과 같을 때, 다음 중 옳은 것은?\end{minipage}\hfill\begin{minipage}[t]{0.4\linewidth}\vspace{0pt}\centering \resizebox{\ifdim\width>\linewidth\linewidth\else\width\fi}{!}{% 91-01(91쪽) — 다항함수 $f(x)$ 의 도함수 $y=f'(x)$ 의 그래프(왼쪽 위에서 내려와 $x=-3$ 에서 $x$축을 지나 극소, $x=-1$ 에서 다시 $x$축을 지나 올라가 $x=1$ 에서 최대, $x=3$ 에서 내려가며 $x$축을 지남). 스캔 화질이 낮아 TikZ 로 다시 그림.
% 원문에 있는 것만 옮긴다: 곡선 · $x=1$ 의 안내 점선 · 라벨 $-3$, $-1$, O, 1, 3, x, y, y=f'(x). 함숫값 눈금은 원문에 없으므로 적지 않는다.
\begin{tikzpicture}[x=0.5cm, y=0.5cm, line width=0.6pt]
  \draw[->, >=stealth, line width=0.7pt] (-4.2,0) -- (4.4,0) node[below=1pt, font=\small] {$x$};
  \draw[->, >=stealth, line width=0.7pt] (0,-2.9) -- (0,5.1) node[left=1pt, font=\small] {$y$};
  \draw[densely dotted, line width=0.4pt] (1,0) -- (1,3.9);
  \draw plot[smooth, tension=0.75] coordinates {(-3.7,3.5) (-3.45,2.4) (-3.2,1.1) (-3,0) (-2.75,-0.95) (-2.5,-1.55) (-2.3,-1.85) (-2.1,-1.9) (-1.9,-1.8) (-1.65,-1.5) (-1.4,-1) (-1.2,-0.5) (-1,0) (-0.7,0.85) (-0.35,1.7) (0,2.4) (0.35,3.1) (0.7,3.65) (1,3.9) (1.35,3.8) (1.7,3.45) (2.05,2.9) (2.4,2.1) (2.7,1.2) (3,0) (3.2,-0.85) (3.4,-1.8) (3.55,-2.45)};
  \node[below=1pt, font=\small] at (-3,0) {$-3$};
  \node[above left=-2pt and -5pt, font=\small] at (-1,0) {$-1$};
  \node[below left=-1pt and -4pt, font=\small] at (0,0) {$\pt{O}$};
  \node[below=1pt, font=\small] at (1,0) {$1$};
  \node[above right=-2pt and -5pt, font=\small] at (3,0) {$3$};
  \node[anchor=west, font=\small] at (0.35,4.45) {$y=f'(x)$};
\end{tikzpicture}
}\end{minipage}\par
\par\medskip\par\noindent\hangindent=1.4em\hangafter=1 {\small ①}\ $f(x)$는 구간 $(-\infty{,}\mkern3mu\mathopen{}-3)$에서 감소한다.\par\noindent\hangindent=1.4em\hangafter=1 {\small ②}\ $f(x)$는 구간 $(-3{,}\mkern3mu\mathopen{}-1)$에서 증가한다.\par\noindent\hangindent=1.4em\hangafter=1 {\small ③}\ $f(x)$는 구간 $(-1{,}\mkern3mu\mathopen{}1)$에서 증가한다.\par\noindent\hangindent=1.4em\hangafter=1 {\small ④}\ $f(x)$는 구간 $(1{,}\mkern3mu\mathopen{}3)$에서 감소한다.\par\noindent\hangindent=1.4em\hangafter=1 {\small ⑤}\ $f(x)$는 구간 $(3{,}\mkern3mu\mathopen{}\infty)$에서 증가한다.\par\medskip
